%-----------------------------------------------------------------------
\section{Background and Mathematical Foundations}
\label{sec:background}
%-----------------------------------------------------------------------

To comprehend the distinct performance divergences intrinsic to post-quantum encryption and authentication, we survey the core mathematical assumptions underpinning the NIST algorithm portfolio.

\subsection{Mathematical Foundations of PQC}

\subsubsection{Lattice-Based Cryptography (Module-LWE)}
The security of lattice-based mechanisms anchors on Learning With Errors (LWE) and Module-LWE \cite{regev2009lattices}. Given a public matrix $\mathbf{A} \in \mathbb{Z}_q^{k \times k}$ and a secret vector $\mathbf{s}$, an adversary is challenged to find $\mathbf{s}$ given $\mathbf{t} = \mathbf{A}\mathbf{s} + \mathbf{e} \pmod q$, where $\mathbf{e}$ is a small error polynomial. 

\textbf{ML-KEM (FIPS 203)} operates over the ring $R_q = \mathbb{Z}_q[X]/(X^{256} + 1)$ with $q = 3329$. It provides efficient key encapsulation but demands substantial register availability for polynomial multiplications. 

\textbf{ML-DSA (FIPS 204)} is similar to ML-KEM, except it utilizes the "Fiat Shamir with Aborts" method. It rejects random coefficients when they go over certain limits while producing a signature. It starts over to ensure secret key leakage is not compromised. This non-deterministic algorithm can consume considerable power for Class 1 devices, especially when the CPU cache is limited. Moreover, when key sizes go over 2.5 KB, they become fragmented.

\subsubsection{\ntru{} Lattices and FN-DSA}
\textbf{FN-DSA (Draft FIPS 206)}\footnote{NIST FIPS 206 remains a draft standard pending finalization of specifications.} overcomes the issue of signature size, which is a limitation in ML-DSA, where the signature is reduced to 666 bytes to achieve 128-bit security strength. This is done through dense lattices \cite{hoffstein1998ntru}. However, the signature generation is continuous, involving Gaussian sampling in high-precision floating-point arithmetic. 

\subsubsection{Hash-Based Signatures (SLH-DSA)}
\textbf{SLH-DSA (FIPS 205)} secures digital signatures upon proven hash function properties (e.g., SHA-256), circumventing unproven algebraic constructs \cite{bernstein2019sphincs}. It constructs a hyper-tree combining "Few-Time Signatures" (FORS) and "eXtended Merkle Signature Schemes" \cite{bernstein2019sphincs}.

\subsubsection{Code-Based Cryptography (HQC)}
\textbf{HQC} derives security from the syndrome decoding problem \cite{nist2025hqc}. It bypasses NTT convolutions entirely, opting for simpler vector additions over finite fields GF(2), offering ultra-low power execution despite large public keys.

\begin{table}[htbp]
\caption{Cryptographic Primitive Size Comparison: Classical vs. PQC Baseline}
\label{tab:classical_vs_pqc}
\centering
\begin{tabular}{llcc}
\toprule
\textbf{Category} & \textbf{Algorithm} & \textbf{Public Key (B)} & \textbf{Ciphertext/Sig (B)} \\
\midrule
\multirow{2}{*}{Classical} 
& RSA-3072 & 384 & 384 \\
& ECDSA (P-256) & 64 & 64 \\
\midrule
\multirow{2}{*}{PQC KEM} 
& \kyber{}-512 & 800 & 768 \\
& \kyber{}-768 & 1,184 & 1,088 \\
\midrule
\multirow{4}{*}{PQC Signature} 
& \dilithium{}-2 & 1,312 & 2,420 \\
& \dilithium{}-3 & 1,952 & 3,293 \\
& \falcon{}-512 & 897 & 666 \\
& \sphincs{}-128f & 32 & 17,088 \\
\bottomrule
\end{tabular}
\end{table}

\subsection{Defined Deployment Targets: RFC 7228}

To maintain rigid boundaries reflective of genuine embedded hardware, we evaluate using the strict taxonomy formalized by the IETF's RFC 7228 \cite{bormann2014coap}:

\begin{itemize}
 \item \textbf{Class 0 (C0):} Devices exhibiting $<10$ KB RAM and $<100$ KB Flash storage. Commonly bare-metal sensory motes devoid of standard TCP/IP networking, necessitating proxy gateways.
 \item \textbf{Class 1 (C1):} Devices exhibiting $\approx 10-50$ KB RAM and $\approx 100-250$ KB Flash. Suitable for lightweight Datagram TLS (DTLS) and UDP implementations (e.g., nRF52840 nodes).
 \item \textbf{Class 2 (C2):} Devices exhibiting $\approx 50-250$ KB RAM and $\approx 250$ KB - 1 MB Flash. Readily capable of implementing nearly full TCP/IPv6 stacks (e.g., robust ESP32 environments).
\end{itemize}

\subsection{Related Work}

Embedded PQC benchmarking now covers a broad range of platforms, but each study addresses only a slice of the problem. The pqm4 project \cite{kannwischer2019pqm4} established cycle-accurate instruction counts for ARM Cortex-M4 — still the most widely cited baseline for constrained crypto. Fitzgibbon and Ottaviani \cite{fitzgibbon2024benchmarking} moved the evaluation to Raspberry Pi 4 under Linux, finding KEM operations practical but signature schemes memory-bound. Lopez et al. \cite{lopez2025post} isolated a subtler issue: garbage collection stalls that appear only when the radio is active during a PQC handshake on ESP32-C3. Silva et al. \cite{secrypt2025esp32} tested ML-DSA and FN-DSA on ESP32 under FreeRTOS at SECRYPT 2025 and documented severe FPU emulation penalties for FALCON's Gaussian sampler. Demir et al. \cite{demir2025power} took a different angle entirely, showing through power trace analysis that bolting constant-time protections onto ML-KEM can triple energy draw.

What none of these studies attempted is a single evaluation covering all five algorithm families across all three RFC 7228 device classes, with energy modeling and multi-criteria ranking included. That gap is what this paper fills.
